\section{Introduction}\label{sec:intro}
Condition-based maintenance and safety-critical operation both rest on prognostics: an estimate of
an asset's remaining useful life (RUL) that a planner can act on. What a maintenance or replacement
decision actually needs is not a point estimate but a trustworthy \emph{lower bound} on remaining
life, with a stated confidence that the true life will not fall below it. An interval that is too
optimistic invites in-service failure, with the attendant safety and cost consequences; one that is
too conservative forces premature replacement and lost availability. The usefulness of a prognostic
interval is therefore set by whether its stated confidence is honest.

Conformal prediction has become an attractive way to attach such a guarantee to a data-driven RUL
model, because it certifies a finite-sample, distribution-free coverage level without assuming a
noise model~\cite{vovk2005algorithmic,papadopoulos2002inductive,lei2018distribution,angelopoulos2023gentle,javanmardi2023conformal}.
That guarantee, however, holds only when the calibration and deployment data are exchangeable, and
this is precisely the assumption that fails in the situation reliability engineering meets most
often: a model calibrated on assets run under one operating regime (load, speed, temperature) is
deployed on assets run under another. The degradation distribution shifts, exchangeability is lost,
and the interval miscovers with no outward sign that its confidence is no longer valid.

The two established ways to repair coverage under such a shift do not fit the reliability setting.
Reweighting methods, namely covariate-shift conformal prediction~\cite{tibshirani2019conformal} and
its non-exchangeable generalisation~\cite{barber2023conformal}, restore coverage, but only by
estimating the likelihood ratio between the target and source operating distributions, which
requires representative failure data from the target regime. A newly commissioned asset, or a regime
with no run-to-failure history, does not have that data. Dimensional-analysis transfer, in which a
model is expressed in dimensionless variables so that it generalises across scales and
systems~\cite{oppenheimer2023multiscale,therrien2024buckingham,girard2024dimensionless}, needs no
target data, but it improves only the \emph{point} prediction and carries no coverage guarantee.
One family has the guarantee but needs data the operator lacks; the other needs no data but offers
no guarantee.

This paper closes that gap with a deployable method, Similarity-Calibrated Conformal (SCC)
prediction, summarised in Fig.~\ref{fig:overview}. SCC computes the conformal calibration in a
dimensionless coordinate chosen, from the governing physics, so that the score distribution is
invariant across operating regimes under dynamic similitude. Coverage then transfers with no target
failure data, and the residual shortfall when similitude is imperfect is bounded by a physical
quantity computable before deployment. A companion runtime check tells the operator when that
premise is supported, so the method declines to certify rather than miscover silently.

The contributions are practical capabilities for reliable prognostics under operating-regime
transfer:
\begin{enumerate}
\item a calibration step in dimensionless ($\Pi$) coordinates that restores conformal coverage
across operating regimes without target-regime failure data (Section~\ref{sec:method});
\item an a-priori coverage bound computable from physical parameters alone: the cross-regime
coverage shortfall is at most $2L\lVert\Delta\boldsymbol\psi\rVert$ in a similitude-departure
quantity $\boldsymbol\psi$, so the worst-case loss can be checked before an asset is deployed
(Section~\ref{sec:method}; derivation in~\ref{app:proof});
\item a runtime similitude diagnostic returning holds, violated, or indeterminate, which gives SCC a
self-aware operating envelope and a conservative fallback when its premise is unsupported
(Section~\ref{sec:method});
\item a non-circular validation on a controlled process-reliability testbed, stress-testing the
guarantee against a graded, physically meaningful departure from similitude
(Sections~\ref{sec:design}--\ref{sec:results}).
\end{enumerate}

We are deliberate about scope. The guarantee is conditional on two physical hypotheses that are
falsifiable rather than assumed; the validation uses a controlled simulation in which similitude is
real but is never assumed by the estimator; and the one field benchmark (FEMTO bearings) is
correctly reported by the diagnostic as outside the method's envelope rather than forced into a
positive result. Validation on adequately powered field data is the natural next step.
